% Appendix: Technical Details
\appendix
\section{Technical Details and Derivations}
\label{app:technical_details}

\subsection{Clifford Algebra Conventions}

\subsubsection{Signature and Basis}

We use the $(-, +, +, +, ...)$ metric signature. The Clifford algebra $Cl(1,9)$ has generators $\Gamma^A$ satisfying:

\begin{equation}
    \{\Gamma^A, \Gamma^B\} = 2\eta^{AB} \cdot \mathbf{1}_{32}
\end{equation}

where $A, B \in \{0, 1, ..., 9\}$ and $\eta^{AB} = \text{diag}(-1, +1, ..., +1)$.

The 10D generators can be decomposed:
\begin{align}
    \Gamma^\mu &= \gamma^\mu \otimes \mathbf{1}_2 \otimes \mathbf{1}_2 \quad (\mu = 0,1,2,3) \\
    \Gamma^{4+a} &= \gamma^5 \otimes \tau^a \otimes \mathbf{1}_2 \quad (a = 1,2) \\
    \Gamma^{7+i} &= \gamma^5 \otimes \tau^3 \otimes \lambda^i \quad (i = 1,...,3)
\end{align}

where $\gamma^\mu$ are 4D Dirac matrices, $\tau^a$ are Pauli matrices, and $\lambda^i$ are Gell-Mann matrices.

\subsubsection{Charge Conjugation Matrix}

The charge conjugation matrix $C_{10}$ satisfies:
\begin{equation}
    C_{10}^{-1} \Gamma^A C_{10} = -(\Gamma^A)^T
\end{equation}

Explicitly:
\begin{equation}
    C_{10} = C_4 \otimes i\sigma^2 \otimes i\sigma^2
\end{equation}

where $C_4 = i\gamma^2\gamma^0$ in the Dirac representation.

\subsection{Torsion Tensor Components}

\subsubsection{Full Torsion Decomposition}

The torsion tensor $T^\lambda_{\mu\nu} = -T^\lambda_{\nu\mu}$ decomposes into three irreducible parts:

\begin{equation}
    T_{\lambda\mu\nu} = \frac{2}{3}(t_{\lambda\mu\nu} - t_{\lambda\nu\mu}) + \frac{1}{3}(g_{\lambda\mu}V_\nu - g_{\lambda\nu}V_\mu) + \epsilon_{\lambda\mu\nu\rho}A^\rho
\end{equation}

where:
\begin{itemize}
    \item $t_{\lambda\mu\nu}$: Trace-free tensor part ($t^\lambda_{\lambda\nu} = 0$)
    \item $V_\mu = T^\lambda_{\lambda\mu}$: Torsion vector (trace)
    \item $A^\rho = \frac{1}{6}\epsilon^{\rho\lambda\mu\nu}T_{\lambda\mu\nu}$: Torsion axial vector
\end{itemize}

\subsubsection{CTUFT Torsion Ansatz}

In CTUFT, we make the specific ansatz:
\begin{equation}
    T_{\lambda\mu\nu} = \tau_0 \cdot f(E/E_c) \cdot \epsilon_{\lambda\mu\nu\rho}u^\rho
\end{equation}

where $u^\rho$ is a unit timelike vector and $f(x)$ is the spectral function:
\begin{equation}
    f(x) = \frac{1}{1 + x^2} \quad \text{with} \quad x = E/E_c
\end{equation}

\subsection{Spectral Dimension Calculation}

\subsubsection{Heat Kernel Method}

The spectral dimension is defined via the return probability $P(\sigma)$ of a diffusion process:

\begin{equation}
    d_s = -2\frac{d\ln P(\sigma)}{d\ln \sigma}
\end{equation}

For a fractal spacetime, the heat kernel trace is:
\begin{equation}
    K(t) = \Tr(e^{t\Delta}) = \frac{1}{(4\pi t)^{d_s/2}} \cdot V_{eff}(t)
\end{equation}

\subsubsection{Running Spectral Dimension}

In CTUFT, the spectral dimension runs as:

\begin{equation}
    d_s(E) = 4 + \frac{6}{1 + (E/E_c)^2}
\end{equation}

Derivation:
\begin{enumerate}
    \item At high energy ($E \gg E_c$): $d_s \to 10$
    \item At low energy ($E \ll E_c$): $d_s \to 4$
    \item Transition occurs at $E \sim E_c = 2 \times 10^{21}$ GeV
\end{enumerate}

\subsection{Field Equations in Detail}

\subsubsection{Einstein-Cartan Equations}

With torsion, the field equations generalize to:

\begin{equation}
    G_{\mu\nu} + \Lambda g_{\mu\nu} = \kappa^2 T_{\mu\nu}^{(total)}
\end{equation}

where the total energy-momentum includes torsion contributions:

\begin{equation}
    T_{\mu\nu}^{(total)} = T_{\mu\nu}^{(matter)} + T_{\mu\nu}^{(torsion)}
\end{equation}

The torsion energy-momentum tensor:
\begin{equation}
    T_{\mu\nu}^{(torsion)} = \frac{1}{\kappa^2}\left(2Q_{\mu\lambda\rho}Q_{\nu}^{\lambda\rho} - Q_{\lambda\rho\sigma}Q^{\lambda\rho\sigma}g_{\mu\nu}\right)
\end{equation}

where $Q_{\mu\nu\lambda} = -\frac{1}{2}T_{\mu\nu\lambda}$ is the contortion tensor.

\subsubsection{Torsion Field Equation}

The torsion field obeys:

\begin{equation}
    \nabla_\lambda Q^{\lambda\mu\nu} - \frac{1}{2}T^\lambda_{\lambda\rho}Q^{\rho\mu\nu} = \frac{\kappa^2}{2}S^{\mu\nu}
\end{equation}

where $S^{\mu\nu}$ is the spin density of matter.

For the CTUFT ansatz with no matter spin:
\begin{equation}
    \Box \tau + m_\tau^2 \tau = 0
\end{equation}

where $m_\tau = \tau_0^{-1} \approx 2 \times 10^{18}$ GeV is the torsion mass scale.

\subsection{Gauge Field Equations}

\subsubsection{Electromagnetic Sector}

Maxwell equations with torsion corrections:

\begin{align}
    \nabla_\mu F^{\mu\nu} &= J^\nu \\
    \partial_{[\mu}F_{\nu\lambda]} &= 0
\end{align}

where the covariant derivative includes torsion:
\begin{equation}
    \nabla_\mu F^{\mu\nu} = \partial_\mu F^{\mu\nu} + \Gamma^\mu_{\mu\lambda}F^{\lambda\nu} + \Gamma^\nu_{\mu\lambda}F^{\mu\lambda}
\end{equation}

Torsion-induced effective current:
\begin{equation}
    J^\nu_{(torsion)} = Q_{\mu\lambda}^{\nu}F^{\mu\lambda} - Q_{\mu\lambda}^{\mu}F^{\lambda\nu}
\end{equation}

\subsubsection{Weak Interaction Sector}

The $SU(2)_L$ gauge fields $W_\mu^a$ satisfy:

\begin{equation}
    D_\mu W^{\mu\nu,a} - ig\epsilon^{abc}W_\mu^b W^{\mu\nu,c} = g J^{\nu,a}
\end{equation}

where $D_\mu$ includes both gauge and gravitational (torsion) connections.

\subsubsection{Strong Interaction Sector}

$SU(3)_C$ Yang-Mills equations:

\begin{equation}
    D_\mu G^{\mu\nu,a} = g_s J^{\nu,a}
\end{equation}

with field strength:
\begin{equation}
    G^{\mu\nu,a} = \partial^\mu A^{\nu,a} - \partial^\nu A^{\mu,a} + g_s f^{abc}A^{\mu,b}A^{\nu,c}
\end{equation}

\subsection{Fermion Mass Matrix}

\subsubsection{Yukawa Couplings from Geometry}

Fermion masses arise from torsion-induced interactions:

\begin{equation}
    \mathcal{L}_{Yukawa} = -\sum_{i,j} Y_{ij}^{(f)} \bar{\psi}_{L,i}^{(f)} \phi \psi_{R,j}^{(f)} + \text{h.c.}
\end{equation}

The Yukawa matrices are determined geometrically:
\begin{equation}
    Y_{ij}^{(f)} = \tau_0 \cdot \langle\psi_i^{(f)}|\hat{T}|\psi_j^{(f)}\rangle
\end{equation}

where $\hat{T}$ is the torsion operator.

\subsubsection{Mass Hierarchy Formula}

The mass ratios follow:

\begin{equation}
    \frac{m_{f_i}}{m_{f_j}} = \exp\left(-\frac{\Delta n_{ij}}{n_0}\right) \cdot \left(\frac{\tau_0}{\tau_{ij}}\right)^{\alpha_f}
\end{equation}

where:
\begin{itemize}
    \item $\Delta n_{ij}$: Difference in generation quantum number
    \item $n_0 \approx 1.8$: Characteristic scale
    \item $\tau_{ij}$: Family-specific torsion coupling
    \item $\alpha_f = 1$ (quarks), $1/2$ (leptons)
\end{itemize}

\subsubsection{Specific Predictions}

\begin{align}
    \frac{m_c}{m_t} &= 1.5 \times 10^{-3} \quad \text{(pred)} \quad vs \quad 1.5 \times 10^{-3} \quad \text{(obs)} \\
    \frac{m_\mu}{m_\tau} &= 5.9 \times 10^{-2} \quad \text{(pred)} \quad vs \quad 5.9 \times 10^{-2} \quad \text{(obs)} \\
    \frac{m_s}{m_b} &= 2.3 \times 10^{-2} \quad \text{(pred)} \quad vs \quad 2.4 \times 10^{-2} \quad \text{(obs)}
\end{align}

\subsection{Renormalization Group Equations}

\subsubsection{Running Couplings}

The gauge couplings run with energy as:

\begin{equation}
    \frac{1}{\alpha_i(E)} = \frac{1}{\alpha_i(E_0)} - \frac{b_i}{2\pi}\ln\left(\frac{E}{E_0}\right) + \delta_i(E)
\end{equation}

where $b_i$ are the standard one-loop coefficients and $\delta_i(E)$ are CTUFT corrections:

\begin{equation}
    \delta_i(E) = c_i \cdot \frac{E^2}{E_c^2} \cdot \ln\left(\frac{E_c}{E}\right)
\end{equation}

\subsubsection{Unification Conditions}

At $E = E_c$, all couplings unify:
\begin{equation}
    \alpha_1(E_c) = \alpha_2(E_c) = \alpha_3(E_c) = \alpha_{GUT}
\end{equation}

CTUFT prediction:
\begin{equation}
    \alpha_{GUT}^{-1} = 26.3 \pm 0.5 \quad \text{(at 2-sigma)}
\end{equation}

\subsection{Numerical Values of Key Parameters}

\begin{table}[htbp]
\centering
\caption{Fundamental Parameters in CTUFT}
\begin{tabular}{@{}llll@{}}
\toprule
\textbf{Parameter} & \textbf{Symbol} & \textbf{Value} & \textbf{Units} \\
\midrule
Fundamental time scale & $\tau_0$ & $5 \times 10^{-3}$ & GeV$^{-1}$ \\
Fundamental energy scale & $E_c$ & $2 \times 10^{21}$ & GeV \\
Fundamental length scale & $\ell_c$ & $10^{-35}$ & m \\
Planck mass & $M_{Pl}$ & $1.22 \times 10^{19}$ & GeV \\
Reduced Planck mass & $m_P$ & $2.44 \times 10^{18}$ & GeV \\
Gravitational coupling & $\kappa$ & $4.1 \times 10^{-19}$ & GeV$^{-1}$ \\
Cosmological constant & $\Lambda$ & $1.1 \times 10^{-52}$ & m$^{-2}$ \\
Hubble constant & $H_0$ & $70$ & km/s/Mpc \\
\bottomrule
\end{tabular}
\label{tab:fundamental_params}
\end{table}

\begin{table}[htbp]
\centering
\caption{Derived Parameters and Their Values}
\begin{tabular}{@{}llll@{}}
\toprule
\textbf{Parameter} & \textbf{Formula} & \textbf{Value} & \textbf{Units} \\
\midrule
Torsion mass scale & $m_\tau = 1/\tau_0$ & $2 \times 10^{18}$ & GeV \\
Unification scale & $E_{GUT} = E_c$ & $2 \times 10^{21}$ & GeV \\
Proton lifetime & $\tau_p \sim E_c^4/M_{Pl}^5$ & $\gg 10^{35}$ & years \\
PBH mass peak & $M_{PBH} \sim (t_{form}/t_{Pl})^2 M_{Pl}$ & $10^{-12}$ & $M_\odot$ \\
Dark energy density & $\rho_\Lambda = 3\tau_0^2/(2\kappa^2)$ & $5.4 \times 10^{-27}$ & kg/m$^3$ \\
\bottomrule
\end{tabular}
\label{tab:derived_params}
\end{table}

\subsection{Computational Implementation}

\subsubsection{Numerical Methods}

For solving CTUFT field equations numerically:

\begin{enumerate}
    \item \textbf{Spectral methods}: For smooth solutions
    \item \textbf{Finite elements}: For complex geometries
    \item \textbf{Lattice regularization}: For quantum corrections
\end{enumerate}

Example: Spherically symmetric static solution
\begin{equation}
    ds^2 = -e^{2\Phi(r)}dt^2 + e^{2\Lambda(r)}dr^2 + r^2 d\Omega^2
\end{equation}

Field equations become ODEs:
\begin{align}
    \frac{d\Phi}{dr} &= \frac{1 - e^{2\Lambda}}{2r} + \kappa^2 r e^{2\Lambda}(\rho + \tau_{00}) \\
    \frac{d\Lambda}{dr} &= \frac{e^{2\Lambda} - 1}{2r} + \kappa^2 r e^{2\Lambda}(p + \tau_{rr})
\end{align}

\subsubsection{Simulation Parameters}

Recommended numerical values for simulations:
\begin{itemize}
    \item Grid resolution: $\Delta x \sim 10^{-3}\tau_0$ (for torsion effects)
    \item Time step: $\Delta t \sim 10^{-2}\tau_0$ (for stability)
    \item Boundary conditions: Outgoing wave (for radiation)
\end{itemize}

\subsection{References for Technical Details}

For readers seeking deeper technical background:
\begin{itemize}
    \item \textbf{Clifford algebras}: Porteous, ``Clifford Algebras and the Classical Groups''
    \item \textbf{Torsion gravity}: Hehl et al., Rev. Mod. Phys. 48, 393 (1976)
    \item \textbf{Spectral dimension}: Ambjørn et al., Phys. Rev. Lett. 100, 091304 (2008)
    \item \textbf{Quantum gravity}: Rovelli, ``Quantum Gravity'' (Cambridge, 2004)
\end{itemize}